\documentclass[11pt]{article}

\usepackage[T1]{fontenc}
\usepackage{lmodern}
\usepackage{amsmath,amssymb,amsthm,mathtools}
\usepackage{microtype}
\usepackage[margin=0.88in]{geometry}
\usepackage{booktabs}
\usepackage{longtable}
\usepackage{array}
\usepackage{enumitem}
\usepackage{xcolor}
\usepackage{listings}
\usepackage{hyperref}
\usepackage{cleveref}
\usepackage{fancyhdr}
\usepackage{lastpage}

\hypersetup{
  colorlinks=true,
  linkcolor=blue!45!black,
  urlcolor=blue!45!black,
  citecolor=blue!45!black,
  pdftitle={Contender 7: Autonomous Operator Reflection in Lean},
  pdfsubject={A constructive large-number proposal for BusyLean},
  pdfkeywords={Lean, large numbers, reflection, W-types, proof theory}
}

\pagestyle{fancy}
\setlength{\headheight}{14pt}
\fancyhf{}
\lhead{BusyLean Contender 7}
\rhead{Autonomous operator reflection}
\cfoot{\thepage\ of \pageref*{LastPage}}

\lstdefinelanguage{lean}{
  keywords={abbrev,def,theorem,inductive,structure,where,match,with,fun,let,in,
    namespace,end,open,import,by,have,show,exact,apply,intro,change,rw,refine,
    obtain,termination_by,private,example},
  sensitive=true,
  morecomment=[l]{--},
  morecomment=[s]{/-}{-/},
  morestring=[b]"
}
\lstset{
  language=lean,
  basicstyle=\ttfamily\footnotesize,
  keywordstyle=\bfseries\color{blue!45!black},
  commentstyle=\itshape\color{black!58},
  stringstyle=\color{green!35!black},
  columns=fullflexible,
  keepspaces=true,
  breaklines=true,
  breakatwhitespace=false,
  showstringspaces=false,
  frame=single,
  framerule=0.25pt,
  rulecolor=\color{black!25},
  xleftmargin=0.6em,
  xrightmargin=0.6em,
  aboveskip=0.8em,
  belowskip=0.8em,
  literate=
    {→}{{$\to$}}1 {←}{{$\leftarrow$}}1 {ℕ}{{$\mathbb{N}$}}1
    {∀}{{$\forall$}}1 {∃}{{$\exists$}}1 {≤}{{$\le$}}1
    {<}{{$<$}}1 {≥}{{$\ge$}}1 {λ}{{$\lambda$}}1
    {ψ}{{$\psi$}}1 {Ω}{{$\Omega$}}1 {ω}{{$\omega$}}1
    {·}{{$\cdot$}}1 {≈}{{$\approx$}}1 {≃}{{$\simeq$}}1 {∈}{{$\in$}}1
    {⟨}{{$\langle$}}1 {⟩}{{$\rangle$}}1
}

\newtheorem{theorem}{Theorem}[section]
\newtheorem{lemma}[theorem]{Lemma}
\newtheorem{corollary}[theorem]{Corollary}
\theoremstyle{definition}
\newtheorem{definition}[theorem]{Definition}
\newtheorem{construction}[theorem]{Construction}
\theoremstyle{remark}
\newtheorem{remark}[theorem]{Remark}
\newtheorem{status}[theorem]{Validation status}

\newcommand{\N}{\mathbb{N}}
\newcommand{\Fam}{\mathcal{F}}
\newcommand{\Trans}{\mathcal{T}}
\newcommand{\Run}{\operatorname{Run}}
\newcommand{\Lift}{\operatorname{Lift}}
\newcommand{\OpStage}{\operatorname{OpStage}}
\newcommand{\Refl}{\operatorname{Ref}}
\newcommand{\Index}{\mathsf{Index}}
\newcommand{\Pfive}{\mathsf{P5}}
\newcommand{\lean}[1]{\texttt{#1}}
\newcommand{\codepath}[1]{\texttt{#1}}

\title{\textbf{Contender 7: Autonomous Operator Reflection in Lean}\\[0.45em]
\large A constructive large-number proposal for the BusyLean repository}
\author{Research artifact prepared for the BusyLean project}
\date{12 October 2026}

\begin{document}
\maketitle

\begin{abstract}
The BusyLean repository formalizes a Lean edition of the ``name the biggest
number'' contest.  Its current Contender~6 applies forest reflection over a
finite-type W-language along a universal W-tree that joins every tree definable
in that language.  This article proposes a next entry, \emph{Contender~7}, by
lifting the same construction one type level: from oracle families to
\emph{transformers of oracle families}.  The primitive reflection operation is
packaged as a transformer; arbitrary transformers are run along W-trees; the
resulting autonomous lift is itself iterated along the universal tree.  The
closed Lean definition is
\[
  \lean{contender\_7}
  = \OpStage(\Refl,\Index)(\Pfive)(42,42).
\]
The supplied Lean proof source establishes the pointwise inequality
\[
  \texttt{finalFamily }n
  < \texttt{autonomousFamily }n\,n,
\]
and hence
\[
  \texttt{contender\_6}<\texttt{contender\_7}.
\]
The proof uses explicit child occurrences and never assumes monotonicity of
reflection in its oracle, a property known to fail for the underlying object
language.  Under the repository's informal calibration
\(\lambda\approx\psi(\Omega_\omega)\), Contender~6 has scale roughly
\(f_{\lambda^2}(42)\), while the operator-level autonomous closure proposed
here has the heuristic scale \(f_{\lambda^\lambda}(42)\).  The strict
inequality is a formal Lean theorem statement with proof source; the ordinal
calibration is deliberately presented only as a heuristic.  The bundle also
contains source audits, smoke tests, a clean-build script, a unified patch, and
this compiled article.  Because the execution environment used to prepare the
bundle lacks Lean~4.34.0, final kernel compilation and timing remain a required
integration check rather than a claimed result.
\end{abstract}

\tableofcontents

%======================================================================
\section{Problem and repository baseline}

BusyLean changes the usual large-number game in two decisive ways.  A contender
must denote a closed, computable natural number, and superiority over the
previous contender must be justified by a proof checked by Lean's kernel.  The
rules exclude nonconstructive or implementation-only escape hatches such as
\lean{noncomputable}, \lean{partial}, \lean{unsafe}, opaque stand-ins,
\lean{implemented\_by}, external primitives, \lean{native\_decide}, custom
axioms, and appeals to uncomputable functions such as Busy Beaver.  They also
impose separate checking budgets for definitions and proofs.

The repository revision used by this artifact is
\[
\texttt{86b360648163f0249eab79539438b6a83dfe8215}.
\]
At that revision the incumbent is
\begin{lstlisting}
def contender_6 : Nat := stage p5Family index 42
\end{lstlisting}
and the principal comparison theorem is
\begin{lstlisting}
theorem contender_5_lt_contender_6 : contender_5 < contender_6
\end{lstlisting}
The source is organized so that the definition layer can be checked without
normalizing the enormous number.  This proposal preserves that architecture:
definitions live in \codepath{MetaReflection.lean} and
\codepath{Contender7.lean}; proofs live in
\codepath{MetaReflectionFacts.lean} and \codepath{Comparison7.lean}.

\subsection{Why a mere arithmetic wrapper is inadequate}

There are trivial ways to obtain a larger number, for example
\(\lean{contender\_6}+1\) or \(2^{\lean{contender\_6}}\).  Such definitions are
computable and easy to compare, but they violate the scientific intent of the
contest's ``no laziness'' rule: they exploit the previous contender as data and
add only a superficial arithmetic shell.  The new definition must instead
strengthen the \emph{generating mechanism} and should not mention the constant
\lean{contender\_6} at all.

Contender~7 meets this stronger criterion.  Its definition imports the
W-language, universal tree, and ordinary-reflection machinery behind
Contender~6, but not the previous named value.  It introduces a new recursion
at the operator level and proves the old family is embedded strictly below it.

%======================================================================
\section{Contender 6 in functional form}
\label{sec:c6}

The design of Contender~7 is easiest to understand by restating the incumbent
as a functional construction.

\subsection{Families and ordinary reflection}

Let
\[
  \Fam = \N\to\N\to\N
\]
be the type of countable families of unary numerical functions.  In Lean this
is the abbreviation
\begin{lstlisting}
abbrev Family := Nat → Nat → Nat
\end{lstlisting}
The first coordinate selects an oracle member; the second is the numerical
input.  The W-language interpreter is evaluated in an environment whose oldest
variable denotes such a family.  The operation
\(
  \Refl(fs,n)
\)
strictly exceeds both:
\begin{enumerate}[leftmargin=*]
  \item every visible oracle value \(fs(i,n)\) with \(i\le n\), and
  \item every enumerated program of depth at most \(n\), evaluated relative to
        \(fs\).
\end{enumerate}
The implementation is a finite maximum plus a positive offset, so it is
constructive and total.

\subsection{Forest reflection}

The type \lean{WTree Nat} has a leaf constructor and countably branching nodes:
\[
  a ::= \mathsf{leaf}\mid \mathsf{node}(f),\qquad
  f:\N\to \mathsf{WTree}(\N).
\]
Ordinary forest reflection is
\[
\begin{aligned}
  \mathsf{stage}(s,\mathsf{leaf})(n)
    &= \Refl(s,n),\\
  \mathsf{stage}(s,\mathsf{node}(f))(n)
    &= \Refl\bigl((i,m)\mapsto
        \mathsf{stage}(s,f(i))(m),\,n\bigr).
\end{aligned}
\]
At a node, the child stages are retained as a family rather than collapsed
pointwise.  A program may therefore compute which child to call.  This is
crucial: a simple pointwise maximum would erase the dynamic indexing power of
the oracle language.

\subsection{The universal index}

For a program \(t\), \lean{defTree t} interprets it and casts the result to a
\lean{WTree Nat}, using a leaf on type mismatch.  The tree
\lean{treesOfDepth k} joins the trees denoted by every program enumerated at
depth at most \(k\).  Finally,
\begin{lstlisting}
def index : WTree Nat := .node treesOfDepth
\end{lstlisting}
joins those finite-depth joins for every \(k\).  Thus \lean{index} does not
encode a hand-written ordinal notation system; it ranges over every W-tree
that the object language itself defines.

Writing \(\Index\) for this tree and \(s=\Pfive\) for the P5 seed family,
Contender~6 is
\[
  C_6 = \mathsf{stage}(s,\Index)(42).
\]
Equivalently, after packaging reflection as a family transformer, it will be
one autonomous lift of the primitive transformer.  This reformulation reveals
the next natural closure operation.

%======================================================================
\section{From functions to operators}
\label{sec:operators}

\subsection{Transformers}

\begin{definition}[Transformer]
A transformer is a function from oracle families to oracle families:
\[
  \Trans = \Fam\to\Fam.
\]
In Lean:
\begin{lstlisting}
abbrev Transformer := Family → Family
\end{lstlisting}
\end{definition}

The ordinary reflection step is packaged as the transformer
\[
  R(fs)(o,n)=\Refl(fs,n),
\]
where the output coordinate \(o\) is intentionally ignored.
\begin{lstlisting}
def reflectionTransformer : Transformer :=
  fun fs _ n => reflect fs n
\end{lstlisting}
The redundant coordinate is not cosmetic.  Once transformers are combined,
that coordinate becomes the routing channel through which a higher-level
operator selects a child transformer.

\subsection{Running a transformer along a tree}

\begin{definition}[Transformer run]
For \(T\in\Trans\), seed \(s\in\Fam\), and W-tree \(a\), define
\(\Run(T,s,a)\in\Fam\) recursively by
\[
\begin{aligned}
  \Run(T,s,\mathsf{leaf})
    &=T(s),\\
  \Run(T,s,\mathsf{node}(f))
    &=T\bigl((i,n)\mapsto \Run(T,s,f(i))(i,n)\bigr).
\end{aligned}
\]
\end{definition}

The aligned coordinate \(i\) in the recursive clause is important.  Coordinate
\(i\) of the family supplied to \(T\) contains child \(i\), evaluated at child
output coordinate \(i\).  Thus a transformer can inspect or route through the
entire countable forest.

The Lean definition is structural:
\begin{lstlisting}
def runTransformer (op : Transformer) (seed : Family) :
    WTree Nat → Family
  | .leaf => op seed
  | .node f =>
      op (fun i n => runTransformer op seed (f i) i n)
\end{lstlisting}
No general recursion, termination oracle, or postulated well-founded relation
is required; recursion is directly on the inductive W-tree.

\subsection{Autonomous lifting}

\begin{definition}[Autonomous lift]
The lift of a transformer runs it along the universal index:
\[
  \Lift(T)(s)=\Run(T,s,\Index).
\]
\end{definition}

\begin{lstlisting}
def liftTransformer (op : Transformer) : Transformer :=
  fun seed => runTransformer op seed index
\end{lstlisting}
For the primitive reflection transformer, this recovers the incumbent forest
reflection exactly.

\begin{theorem}[Compatibility with Contender 6]
For every seed family \(s\), output coordinate \(o\), and input \(n\),
\[
 \Lift(R)(s)(o,n)=\mathsf{stage}(s,\Index)(n).
\]
\end{theorem}

The supplied Lean theorem is
\begin{lstlisting}
theorem liftReflection_eq_stage (seed : Family) (o n : Nat) :
    liftTransformer reflectionTransformer seed o n =
      stage seed index n
\end{lstlisting}
It follows by structural induction on the W-tree, not by extensional or
ordinal reasoning.

%======================================================================
\section{Iterating the lift itself}
\label{sec:operator-stage}

A single application of \(\Lift\) turns ordinary reflection into the mechanism
behind Contender~6.  Contender~7 makes \(\Lift\) itself the operation that is
iterated.

\subsection{Splicing child transformers}

Given a countable family \((T_i)_{i\in\N}\) of transformers, define their
splice by
\[
  \mathsf{splice}(T_0,T_1,\ldots)(s)(i,n)=T_i(s)(0,n).
\]
Output coordinate \(i\) chooses the \(i\)-th transformer.  Coordinate zero is
the canonical observation point inside the selected child.

\begin{lstlisting}
def spliceTransformers (ops : Nat → Transformer) : Transformer :=
  fun seed i n => ops i seed 0 n
\end{lstlisting}

\subsection{The outer recursion}

\begin{definition}[Operator stage]
For a base transformer \(B\), define a tree-indexed transformer
\(\OpStage(B,a)\) by
\[
\begin{aligned}
  \OpStage(B,\mathsf{leaf})
    &=\Lift(B),\\
  \OpStage(B,\mathsf{node}(f))
    &=\Lift\bigl(\mathsf{splice}
       (i\mapsto\OpStage(B,f(i)))\bigr).
\end{aligned}
\]
\end{definition}

In Lean:
\begin{lstlisting}
def operatorStage (base : Transformer) : WTree Nat → Transformer
  | .leaf => liftTransformer base
  | .node f =>
      liftTransformer
        (spliceTransformers (fun i => operatorStage base (f i)))
\end{lstlisting}
This is the central new definition.  There are now two autonomous W-tree
layers:
\begin{enumerate}[leftmargin=*]
  \item the \emph{inner} universal index inside every application of
        \(\Lift\), and
  \item the \emph{outer} tree controlling how those lifted transformers are
        spliced and lifted again.
\end{enumerate}
Choosing the universal index for the outer tree closes the hierarchy under all
operator stages generated by trees definable in the same W-language.

%======================================================================
\section{The proposed number}
\label{sec:number}

\begin{construction}[Contender 7]
Define
\[
\begin{aligned}
  A &= \OpStage(R,\Index),\\
  F &= A(\Pfive),\\
  C_7 &= F(42,42).
\end{aligned}
\]
\end{construction}

The complete definition layer is only:
\begin{lstlisting}
def autonomousTransformer : Transformer :=
  operatorStage reflectionTransformer index

def autonomousFamily : Family :=
  autonomousTransformer p5Family

def contender_7 : Nat := autonomousFamily 42 42
\end{lstlisting}
The diagonal coordinate \((42,42)\) is deliberate.  The first 42 exposes the
outer transformer's output channel, while the second 42 is the ordinary
numerical input and the visibility bound used by progressivity.  No occurrence
of the constant \lean{contender\_6} appears in these definitions.

\subsection{Computability in principle}

Every component is executable:
\begin{itemize}[leftmargin=*]
  \item program enumerations at a fixed depth are finite lists;
  \item the interpreter of the finite-type W-language is total;
  \item \lean{reflect} computes finite maxima and arithmetic;
  \item \lean{runTransformer} and \lean{operatorStage} recurse structurally on
        inductive W-trees;
  \item \lean{index} is a closed W-tree whose branches are generated by total
        functions; and
  \item the seed \lean{p5Family} is already a total function in the repository.
\end{itemize}
Consequently \lean{contender\_7} is a closed \lean{Nat}.  ``Computable in
principle'' does not mean feasible: evaluating it would require an
inconceivable amount of computation and must not be attempted.  Lean checks
its definition and proofs symbolically, without normalizing the numeral.

%======================================================================
\section{Direct comparison without monotonicity}
\label{sec:comparison}

The key technical risk is monotonicity.  A program in the object language can
inspect an oracle value and branch on it; increasing an oracle pointwise can
therefore decrease a program's output.  Neither \lean{reflect} nor the lifted
operators should be treated as monotone in their oracle argument.  The proof
instead records explicit occurrences.

\subsection{Progressive transformers}

\begin{definition}[Progressivity]
A transformer \(T\) is progressive when
\[
  i\le n \quad\Longrightarrow\quad
  s(i,n)<T(s)(o,n)
\]
for every seed \(s\), every input \(n\), and every output coordinate \(o\).
\end{definition}

The Lean predicate is
\begin{lstlisting}
def Progressive (op : Transformer) : Prop :=
  ∀ (seed : Family) (i o n : Nat),
    i ≤ n → seed i n < op seed o n
\end{lstlisting}
Ordinary reflection is progressive by the repository theorem
\lean{oracle\_lt\_reflect}.

\subsection{Child dominance}

At a node, a visible aligned child is literally one coordinate of the family
supplied to the parent transformer.  Progressivity therefore gives the strict
inequality directly.

\begin{lemma}[Run child]
If \(T\) is progressive and \(i\le n\), then
\[
  \Run(T,s,f(i))(i,n)
  <\Run(T,s,\mathsf{node}(f))(o,n).
\]
\end{lemma}

The corresponding Lean theorem is \lean{runTransformer\_child\_lt}.  No
comparison between two different oracle families is made.

\subsection{A transformer lies below its autonomous lift}

The universal index has the form
\[
  \Index=\mathsf{node}(k\mapsto\mathsf{treesOfDepth}(k)).
\]
At depth zero, the term enumeration is empty.  Every defaulted child of
\lean{treesOfDepth 0} is therefore a leaf.  For any visible coordinate
\(i\le n\), there is an explicit two-edge path
\[
  \mathsf{leaf}
  \longrightarrow \mathsf{treesOfDepth}(0)
  \longrightarrow \Index.
\]
Applying the run-child lemma on each edge yields:

\begin{theorem}[Lift dominance]
If \(T\) is progressive and \(i\le n\), then
\[
  T(s)(i,n)<\Lift(T)(s)(o,n).
\]
\end{theorem}

This is formalized as \lean{transformer\_lt\_lift\_at}.  It immediately gives
that autonomous lifting preserves progressivity.  Splicing also preserves
progressivity when every selected child transformer is progressive.  A
structural induction then proves:

\begin{theorem}[Progressivity of operator stages]
If the base transformer \(B\) is progressive, then
\(\OpStage(B,a)\) is progressive for every W-tree \(a\).
\end{theorem}

A second direct child lemma gives
\[
 \OpStage(B,f(i))(s)(0,n)
 <\OpStage(B,\mathsf{node}(f))(s)(o,n)
 \qquad(i\le n).
\]

\subsection{Contender 6 is a visible outer leaf}

At an outer leaf,
\[
  \OpStage(R,\mathsf{leaf})=\Lift(R).
\]
By the compatibility theorem in \cref{sec:operators}, this is exactly the
Contender~6 family.  The same two explicit outer edges
\[
  \mathsf{leaf}
  \longrightarrow \mathsf{treesOfDepth}(0)
  \longrightarrow \Index
\]
strictly dominate that leaf transformer.  Therefore:

\begin{theorem}[Pointwise formal margin]
For every \(n\in\N\),
\[
  \lean{finalFamily}\ n
  <\lean{autonomousFamily}\ n\ n.
\]
\end{theorem}

The supplied proof is:
\begin{lstlisting}
theorem finalFamily_lt_autonomousFamily (n : Nat) :
    finalFamily n < autonomousFamily n n := by
  have h0 :
      operatorStage reflectionTransformer .leaf p5Family 0 n <
        operatorStage reflectionTransformer
          (treesOfDepth 0) p5Family 0 n := by
    have h := operatorStage_child_lt reflectionTransformer
      reflectionTransformer_progressive p5Family
      (fun i => defTree
        (nthD (BigNumber.allTerms 0) i .zero))
      (i := 0) (o := 0) (n := n) (Nat.zero_le n)
    rw [depthZeroTreeLeaf] at h
    exact h
  have h1 :
      operatorStage reflectionTransformer
          (treesOfDepth 0) p5Family 0 n <
        operatorStage reflectionTransformer
          index p5Family n n :=
    operatorStage_child_lt reflectionTransformer
      reflectionTransformer_progressive p5Family treesOfDepth
      (i := 0) (o := n) (n := n) (Nat.zero_le n)
  have h :
      liftTransformer reflectionTransformer p5Family 0 n <
        operatorStage reflectionTransformer index p5Family n n :=
    Nat.lt_trans h0 h1
  change stage p5Family index n <
    operatorStage reflectionTransformer index p5Family n n
  rw [← liftReflection_eq_stage p5Family 0 n]
  exact h
\end{lstlisting}
Specializing to 42 proves the contest comparison:
\begin{lstlisting}
theorem contender_6_lt_contender_7 :
    contender_6 < contender_7 :=
  finalFamily_lt_autonomousFamily 42
\end{lstlisting}
By transitivity, the bundle also derives strict inequalities from the exact
Contender~5 port and from proposal P5's number.

%======================================================================
\section{A stronger closure theorem}
\label{sec:closure}

The definition is designed to exceed more than one named constant.  Let \(t\)
be any enumerated object-language program denoting a W-tree.  If \(t\) appears
at finite enumeration depth \(k\), then \lean{defTree t} is a child of
\lean{treesOfDepth k}, and that tree is a child of \lean{index}.  Reusing direct
child dominance at the operator level yields the theorem
\begin{lstlisting}
theorem defined_operator_eventually_lt {k : Nat} {t : Term}
    (ht : t ∈ BigNumber.allTerms k) :
    ∃ N, ∀ n, N ≤ n →
      operatorStage reflectionTransformer (defTree t)
          p5Family 0 n <
        autonomousFamily n n
\end{lstlisting}
Thus every operator-stage whose outer schedule is denoted by a fixed program
of the W-language is eventually dominated by the autonomous family.  This is
the meaningful structural margin: Contender~7 closes not only over the old
number, but over a whole class of independently generated operator hierarchies.

The theorem is still deliberately modest.  It does not claim that every total
Lean transformer is represented, nor does it quantify over arbitrary semantic
well-orderings.  Its scope is exactly the effective W-language already audited
by BusyLean.

%======================================================================
\section{Informal size analysis}
\label{sec:size}

This section is heuristic.  None of its ordinal identifications is used in the
Lean comparison proof.

\subsection{Repository calibration}

Let \(\lambda\) denote the proof-theoretic scale of the finite-type W-language.
The repository informally calibrates it near
\[
  \lambda\approx\psi(\Omega_\omega).
\]
One reflection step relative to an oracle is treated as adding about
\(\lambda\) to the fast-growing index.  Iteration along a W-tree of height
\(\alpha\) then has the rough effect of right multiplication by \(\alpha\).
Because the universal index ranges over the language's own definable trees,
its effective height is heuristically \(\lambda\).  On that account:
\[
  C_6\quad\text{has scale roughly}\quad f_{\lambda\cdot\lambda}(42)
  =f_{\lambda^2}(42).
\]

\subsection{One type-level lift adds an exponent}

Suppose, schematically, that a progressive transformer has increment scale
\(\lambda^e\).  Running it along the universal index composes that increment
through height \(\lambda\), suggesting
\[
  \lambda^e\cdot\lambda=\lambda^{e+1}.
\]
The primitive reflection transformer starts at exponent one.  Its autonomous
lift---the mechanism of Contender~6---is therefore at exponent two.

For a finite outer tree of rank \(h\), \lean{operatorStage} applies one further
universal lift at each outer level.  The finite skeleton is
\[
\begin{array}{c|c}
\text{outer height} & \text{heuristic increment scale}\\
\hline
0 & \lambda^2\\
1 & \lambda^3\\
2 & \lambda^4\\
\vdots & \vdots\\
h & \lambda^{h+2}.
\end{array}
\]
The included script \codepath{scripts/ordinal\_skeleton.py} regression-tests
this finite bookkeeping model.

\subsection{Closing over the universal outer tree}

The outer tree is again \(\Index\), heuristically cofinal through ranks below
\(\lambda\).  Taking the supremum of the finite and transfinite outer stages
therefore suggests
\[
  \sup_{\alpha<\lambda}\lambda^{\alpha+2}
  =\lambda^\lambda.
\]
Accordingly, the proposed calibration is
\[
  C_7\quad\text{roughly}\quad f_{\lambda^\lambda}(42),
  \qquad \lambda\approx\psi(\Omega_\omega).
\]
This is a substantial mechanism-level jump over \(f_{\lambda^2}(42)\), while
remaining within the same constructive object language.

\begin{remark}[Limits of the calibration]
The equation ``Contender~7 is \(f_{\lambda^\lambda}(42)\)'' is not a theorem of
this work.  A rigorous calibration would require an ordinal notation system,
a semantics relating transformer runs to that system, upper and lower bounds,
and a proof that the universal tree enumeration is cofinal in the required
sense.  The present Lean development proves strict numerical inequalities and
structural domination theorems, not ordinal equivalences.
\end{remark}

%======================================================================
\section{Why not claim the System F jump yet?}
\label{sec:systemf}

The existing BusyLean article identifies impredicative polymorphism as a
promising qualitative next step.  A fully satisfactory System~F contender
would require considerably more infrastructure:
\begin{enumerate}[leftmargin=*]
  \item syntax and decidable typing for a suitable presentation of System~F;
  \item substitution, reduction, and a deterministic normalizer;
  \item a formal strong-normalization proof, most naturally through
        reducibility candidates in Lean's impredicative \lean{Prop};
  \item a verified decoder for Church numerals;
  \item a finite bounded syntax enumeration; and
  \item an oracle extension whose normalization proof remains constructive.
\end{enumerate}
Such a development could plausibly dwarf the predicative W-language, but an
uncompiled sketch would not satisfy the contest.  The present proposal instead
makes a smaller foundational change that can be expressed entirely using the
already audited W-language and structural recursion.  It is intended as a
credible, reviewable Contender~7 rather than a promise of a stronger but
unfinished system.

%======================================================================
\section{Lean engineering and dependency discipline}
\label{sec:engineering}

\subsection{Definition/proof separation}

The new files follow the incumbent's time-budget architecture:
\begin{center}
\begin{tabular}{p{0.32\linewidth}p{0.58\linewidth}}
\toprule
File & Role\\
\midrule
\codepath{MetaReflection.lean}
  & Definition-only operator machinery.\\
\codepath{Contender7.lean}
  & Closed family and number definitions.\\
\codepath{MetaReflectionFacts.lean}
  & Progressivity, compatibility, and direct child inequalities.\\
\codepath{Comparison7.lean}
  & Named-contender comparisons and autonomous closure theorem.\\
\codepath{MetaReflectionSmoke.lean}
  & Tiny executable recursion checks.\\
\codepath{Contender7Audit.lean}
  & Axiom reports and theorem signature checks.\\
\bottomrule
\end{tabular}
\end{center}

\subsection{Excluded constructs}

A source audit checks the Lean code, after removing comments, for the forbidden
patterns
\begin{center}
\lean{noncomputable}, \lean{partial}, \lean{unsafe}, \lean{opaque},
\lean{implemented\_by}, \lean{extern}, \lean{native\_decide},
\lean{sorry}, \lean{admit}, and custom \lean{axiom} declarations.
\end{center}
It also checks that the definition layer does not mention
\lean{contender\_6}; that expected declarations exist; and that the aggregate
\codepath{BusyLean.lean} imports all new modules.  The Lean-side audit goes
further: it reconstructs the transitive declaration closure from the kernel
environment and is set to reject dependencies on the named Contender~6 and
Contender~5 constants or selectors.  That closure check becomes authoritative
only when the supplied test is run by Lean.

\subsection{Avoiding accidental normalization}

The comparison proof first proves a general pointwise inequality in a symbolic
variable \(n\), and only then specializes to 42.  The proof never unfolds the
astronomical arithmetic computation.  Rewriting is directed through
\lean{liftReflection\_eq\_stage}, so elaboration needs only definitional
unfolding of the construction, not evaluation of the number.

\subsection{Expected axiom profile}

The mathematical proof uses direct natural-number inequalities, structural
induction, function extensionality in the compatibility lemma, and existing
BusyLean results.  On Lean~4.34.0, the repository permits \lean{propext} and
\lean{Quot.sound}; the audit script rejects any reported axiom outside that
set.  The exact reported set must be established by the final kernel run.

%======================================================================
\section{Validation protocol and current status}
\label{sec:validation}

\subsection{Checks completed in this bundle}

The preparation environment performed the following checks:
\begin{enumerate}[leftmargin=*]
  \item source-level audit of all new Lean and test files;
  \item verification that no forbidden constructive-rule token occurs in code;
  \item verification that the definition layer is independent of the named
        constant \lean{contender\_6};
  \item verification that the expected definitions, theorem names, comparison
        chain, and imports are present;
  \item execution of the finite ordinal-exponent skeleton script;
  \item Python bytecode compilation of the validation scripts;
  \item two-pass \LaTeX{} compilation of this article;
  \item PDF metadata inspection, text extraction, and page rendering; and
  \item visual inspection of the rendered pages for clipping and glyph errors.
\end{enumerate}

\subsection{Decisive check still required}

\begin{status}
The current execution environment contains no Lean or Lake executable and has
no network access from which to install the pinned Lean~4.34.0 binary.  I
therefore do \emph{not} claim that the supplied Lean source has been accepted by
Lean's kernel or that it satisfies the 15-second/60-second budgets on this
machine.  The source has been manually and statically reviewed, but kernel
checking is authoritative.
\end{status}

The bundle supplies \codepath{scripts/check\_contender7.py}.  After applying
the overlay to the recorded BusyLean revision, run
\begin{lstlisting}[language={}]
python3 scripts/check_contender7.py
\end{lstlisting}
The script deletes build directories, builds \lean{BusyLean.Contender7} under
the definition budget, builds the proof layer under the proof budget, executes
the smoke test, runs the axiom audit, and rejects any axiom outside
\(\{\lean{propext},\lean{Quot.sound}\}\).

\subsection{Likely integration failure modes}

The most plausible issues are engineering rather than mathematical:
\begin{itemize}[leftmargin=*]
  \item Lean may require a small elaboration adjustment in a higher-order
        structural recursion equation;
  \item a rewrite may need an explicit unfolding directive to prevent the
        elaborator from choosing an undesirable normal form;
  \item the toolchain may attribute \lean{propext} or \lean{Quot.sound} through
        core lemmas, as already happens elsewhere in the repository; or
  \item clean-build time may exceed the contest budget despite short source.
\end{itemize}
None should be concealed.  If the kernel rejects a term, the theorem is not yet
established, regardless of the paper proof.  The validation report in the
bundle records this distinction.

%======================================================================
\section{Integration instructions}
\label{sec:integration}

The archive contains a \codepath{repo\_overlay/} directory mirroring the
repository root and a unified patch in
\codepath{patches/contender7.patch}.  From a checkout at the recorded commit:
\begin{lstlisting}[language={}]
git apply patches/contender7.patch
python3 scripts/check_contender7.py
\end{lstlisting}
Alternatively copy the overlay into the checkout, preserving paths.  The
aggregate file \codepath{BusyLean.lean} gains four imports:
\begin{lstlisting}
import BusyLean.MetaReflection
import BusyLean.MetaReflectionFacts
import BusyLean.Contender7
import BusyLean.Comparison7
\end{lstlisting}
No existing definition is modified.  The proposal is therefore easy to remove
or review in isolation.

%======================================================================
\section{Further research questions}
\label{sec:further}

\begin{enumerate}[leftmargin=*]
  \item \textbf{Kernel stabilization.}  Compile on Lean~4.34.0, minimize any
        elaboration fixes, and preserve the definition/proof timing split.

  \item \textbf{Richer dependency-closure reporting.}  After the bundled
        rejection test passes, list and classify all inherited constants under
        the old-language and proposal namespaces, as the Contender~6 audit does.

  \item \textbf{Tighter structural theorem.}  Formalize domination of every
        finite iteration of \lean{liftTransformer}, perhaps using canonical
        spine trees and an explicit embedding into \lean{index}.

  \item \textbf{Representation theorem.}  Characterize which progressive
        transformers can be represented by W-language programs or by
        \lean{operatorStage} schedules.

  \item \textbf{Ordinal lower bound.}  Define an ordinal notation semantics for
        the operator hierarchy and prove a lower bound approaching
        \(\lambda^\lambda\), rather than relying on the present calibration
        heuristic.

  \item \textbf{Ordinal upper bound.}  Show that the recursion does not silently
        exceed the proposed calibration because of higher-type oracle routing.
        A matching upper bound would make the ordinal claim meaningful.

  \item \textbf{Higher transformer levels.}  Iterate the construction at types
        \(\Fam\to\Fam\), \((\Fam\to\Fam)\to(\Fam\to\Fam)\), and beyond, while
        retaining a direct, non-monotone comparison principle.

  \item \textbf{Diagonal autonomous tower.}  Encode all finite transformer
        levels in one inductive universe of codes and take a verified diagonal
        across levels.  This may suggest scales analogous to
        \(\lambda\uparrow\uparrow\lambda\), but requires careful typing.

  \item \textbf{Impredicative successor.}  Complete a constructive System~F
        normalization development, then compare a bounded-depth System~F
        maximum against Contender~7 without invoking classical choice.

  \item \textbf{Calculus of Constructions.}  Investigate a Lean-hosted,
        independently normalized dependent type theory in the direction of
        Loader-style numbers, with a faithful and auditable syntax bound.

  \item \textbf{Complexity of checking.}  Measure elaboration and kernel costs
        separately.  A definition may be mathematically compact yet expensive
        for Lean because of generated recursors or closure size.

  \item \textbf{Robustness under toolchain upgrades.}  Test the axiom profile
        and generated code on subsequent Lean releases, documenting changes
        due solely to core implementation details.
\end{enumerate}

%======================================================================
\section{Conclusion}

Contender~7 replaces a one-level autonomous reflection hierarchy by an
autonomous hierarchy of reflection \emph{operators}.  The proposal remains a
closed, structurally recursive Lean natural number and does not mention the
previous contender constant.  Its formal comparison strategy is deliberately
conservative: define progressivity, follow explicit child occurrences, identify
the outer leaf with the Contender~6 family, and compose two strict inequalities.
The same argument proves eventual domination of every operator-stage generated
by a fixed W-language-definable tree.

The construction is therefore substantially more than an arithmetic wrapper.
Under the repository's own proof-theoretic heuristic it moves from
\(\lambda^2\) to \(\lambda^\lambda\), with
\(\lambda\approx\psi(\Omega_\omega)\).  That estimate is not passed off as a
formal theorem.  The formal deliverable is the Lean source for
\lean{contender\_7}, its pointwise comparison theorem, and the associated
closure lemmas.  The remaining gate is unambiguous: apply the patch to the
recorded repository revision and pass the supplied Lean~4.34.0 clean-build,
timing, smoke, and axiom checks.

%======================================================================
\appendix
\section{Definition source}
\label{app:definitions}

\subsection{\codepath{MetaReflection.lean}}
\lstinputlisting[basicstyle=\ttfamily\scriptsize]{../repo_overlay/BusyLean/MetaReflection.lean}

\subsection{\codepath{Contender7.lean}}
\lstinputlisting[basicstyle=\ttfamily\scriptsize]{../repo_overlay/BusyLean/Contender7.lean}

\section{Proof source}
\label{app:proofs}

\subsection{\codepath{MetaReflectionFacts.lean}}
\lstinputlisting[basicstyle=\ttfamily\scriptsize]{../repo_overlay/BusyLean/MetaReflectionFacts.lean}

\subsection{\codepath{Comparison7.lean}}
\lstinputlisting[basicstyle=\ttfamily\scriptsize]{../repo_overlay/BusyLean/Comparison7.lean}

\section{Tests}
\label{app:tests}

\subsection{Smoke test}
\lstinputlisting[basicstyle=\ttfamily\scriptsize]{../repo_overlay/test/MetaReflectionSmoke.lean}

\subsection{Axiom audit}
\lstinputlisting[basicstyle=\ttfamily\scriptsize]{../repo_overlay/test/Contender7Audit.lean}

%======================================================================
\begin{thebibliography}{9}

\bibitem{BusyLean}
V. Reshetnikov,
\emph{BusyLean: name the biggest number, in Lean},
repository revision
\texttt{86b360648163f0249eab79539438b6a83dfe8215}, October 2026.

\bibitem{Roux}
C. Roux,
\emph{name-the-biggest-number},
original Coq contest repository, with Contenders 4 and 5.

\bibitem{Aaronson}
S. Aaronson,
\emph{Who Can Name the Bigger Number?},
expository essay on large-number contests and definability.

\bibitem{Girard}
J.-Y. Girard, Y. Lafont, and P. Taylor,
\emph{Proofs and Types},
Cambridge Tracts in Theoretical Computer Science, 1989.

\bibitem{Pohlers}
W. Pohlers,
\emph{Proof Theory: The First Step into Impredicativity},
Springer, 2009.

\end{thebibliography}

\end{document}
